\documentclass[11pt,preprint]{elsarticle}

\usepackage{lmodern}
%%%% My spacing
\usepackage{setspace}
\setstretch{1.2}
\DeclareMathSizes{12}{14}{10}{10}

% Wrap around which gives all figures included the [H] command, or places it "here". This can be tedious to code in Rmarkdown.
\usepackage{float}
\let\origfigure\figure
\let\endorigfigure\endfigure
\renewenvironment{figure}[1][2] {
    \expandafter\origfigure\expandafter[H]
} {
    \endorigfigure
}

\let\origtable\table
\let\endorigtable\endtable
\renewenvironment{table}[1][2] {
    \expandafter\origtable\expandafter[H]
} {
    \endorigtable
}


\usepackage{ifxetex,ifluatex}
\usepackage{fixltx2e} % provides \textsubscript
\ifnum 0\ifxetex 1\fi\ifluatex 1\fi=0 % if pdftex
  \usepackage[T1]{fontenc}
  \usepackage[utf8]{inputenc}
\else % if luatex or xelatex
  \ifxetex
    \usepackage{mathspec}
    \usepackage{xltxtra,xunicode}
  \else
    \usepackage{fontspec}
  \fi
  \defaultfontfeatures{Mapping=tex-text,Scale=MatchLowercase}
  \newcommand{\euro}{€}
\fi

\usepackage{amssymb, amsmath, amsthm, amsfonts}

\def\bibsection{\section*{References}} %%% Make "References" appear before bibliography


\usepackage[numbers]{natbib}

\usepackage{longtable}
\usepackage[margin=2.3cm,bottom=2cm,top=2.5cm, includefoot]{geometry}
\usepackage{fancyhdr}
\usepackage[bottom, hang, flushmargin]{footmisc}
\usepackage{graphicx}
\numberwithin{equation}{section}
\numberwithin{figure}{section}
\numberwithin{table}{section}
\setlength{\parindent}{0cm}
\setlength{\parskip}{1.3ex plus 0.5ex minus 0.3ex}
\usepackage{textcomp}
\renewcommand{\headrulewidth}{0.2pt}
\renewcommand{\footrulewidth}{0.3pt}

\usepackage{array}
\newcolumntype{x}[1]{>{\centering\arraybackslash\hspace{0pt}}p{#1}}

%%%%  Remove the "preprint submitted to" part. Don't worry about this either, it just looks better without it:

 \def\tightlist{} % This allows for subbullets!

\usepackage{hyperref}
\hypersetup{breaklinks=true,
            bookmarks=true,
            colorlinks=true,
            citecolor=blue,
            urlcolor=blue,
            linkcolor=blue,
            pdfborder={0 0 0}}


% The following packages allow huxtable to work:
\usepackage{siunitx}
\usepackage{multirow}
\usepackage{hhline}
\usepackage{calc}
\usepackage{tabularx}
\usepackage{booktabs}
\usepackage{caption}


\newenvironment{columns}[1][]{}{}

\newenvironment{column}[1]{\begin{minipage}{#1}\ignorespaces}{%
\end{minipage}
\ifhmode\unskip\fi
\aftergroup\useignorespacesandallpars}

\def\useignorespacesandallpars#1\ignorespaces\fi{%
#1\fi\ignorespacesandallpars}

\makeatletter
\def\ignorespacesandallpars{%
  \@ifnextchar\par
    {\expandafter\ignorespacesandallpars\@gobble}%
    {}%
}
\makeatother


% definitions for citeproc citations
\NewDocumentCommand\citeproctext{}{}
\NewDocumentCommand\citeproc{mm}{%
\href{\#cite.\detokenize{#1}}{#2}\nocite{#1}}

\makeatletter
% allow citations to break across lines
\let\@cite@ofmt\@firstofone
% avoid brackets around text for \cite:
\def\@biblabel#1{}
\def\@cite#1#2{{#1\if@tempswa , #2\fi}}
\makeatother
\newlength{\cslhangindent}
\setlength{\cslhangindent}{1.5em}
\newlength{\csllabelwidth}
\setlength{\csllabelwidth}{3em}
\newenvironment{CSLReferences}[2] % #1 hanging-indent, #2 entry-spacing
{\begin{list}{}{%
	\setlength{\itemindent}{0pt}
	\setlength{\leftmargin}{0pt}
	\setlength{\parsep}{0pt}
	% turn on hanging indent if param 1 is 1
	\ifodd #1
	\setlength{\leftmargin}{\cslhangindent}
	\setlength{\itemindent}{-1\cslhangindent}
	\fi
	% set entry spacing
	\setlength{\itemsep}{#2\baselineskip}}}
{\end{list}}

\usepackage{calc}
\newcommand{\CSLBlock}[1]{\hfill\break\parbox[t]{\linewidth}{\strut\ignorespaces#1\strut}}
\newcommand{\CSLLeftMargin}[1]{\parbox[t]{\csllabelwidth}{\strut#1\strut}}
\newcommand{\CSLRightInline}[1]{\parbox[t]{\linewidth - \csllabelwidth}{\strut#1\strut}}
\newcommand{\CSLIndent}[1]{\hspace{\cslhangindent}#1}


\urlstyle{same}  % don't use monospace font for urls
\setlength{\parindent}{0pt}
\setlength{\parskip}{6pt plus 2pt minus 1pt}
\setlength{\emergencystretch}{3em}  % prevent overfull lines
\setcounter{secnumdepth}{5}

%%% Use protect on footnotes to avoid problems with footnotes in titles
\let\rmarkdownfootnote\footnote%
\def\footnote{\protect\rmarkdownfootnote}
\IfFileExists{upquote.sty}{\usepackage{upquote}}{}

%%% Include extra packages specified by user

%%% Hard setting column skips for reports - this ensures greater consistency and control over the length settings in the document.
%% page layout
%% paragraphs
\setlength{\baselineskip}{12pt plus 0pt minus 0pt}
\setlength{\parskip}{12pt plus 0pt minus 0pt}
\setlength{\parindent}{0pt plus 0pt minus 0pt}
%% floats
\setlength{\floatsep}{12pt plus 0 pt minus 0pt}
\setlength{\textfloatsep}{20pt plus 0pt minus 0pt}
\setlength{\intextsep}{14pt plus 0pt minus 0pt}
\setlength{\dbltextfloatsep}{20pt plus 0pt minus 0pt}
\setlength{\dblfloatsep}{14pt plus 0pt minus 0pt}
%% maths
\setlength{\abovedisplayskip}{12pt plus 0pt minus 0pt}
\setlength{\belowdisplayskip}{12pt plus 0pt minus 0pt}
%% lists
\setlength{\topsep}{10pt plus 0pt minus 0pt}
\setlength{\partopsep}{3pt plus 0pt minus 0pt}
\setlength{\itemsep}{5pt plus 0pt minus 0pt}
\setlength{\labelsep}{8mm plus 0mm minus 0mm}
\setlength{\parsep}{\the\parskip}
\setlength{\listparindent}{\the\parindent}
%% verbatim
\setlength{\fboxsep}{5pt plus 0pt minus 0pt}



\begin{document}



\begin{frontmatter}  %

\title{Netflix Report}

% Set to FALSE if wanting to remove title (for submission)












\vspace{1cm}





\vspace{0.5cm}

\end{frontmatter}

\setcounter{footnote}{0}



%________________________
% Header and Footers
%%%%%%%%%%%%%%%%%%%%%%%%%%%%%%%%%
\pagestyle{fancy}
\chead{}
\rhead{}
\lfoot{}
\rfoot{\footnotesize Page \thepage}
\lhead{}
%\rfoot{\footnotesize Page \thepage } % "e.g. Page 2"
\cfoot{}

%\setlength\headheight{30pt}
%%%%%%%%%%%%%%%%%%%%%%%%%%%%%%%%%
%________________________

\headsep 35pt % So that header does not go over title




\section{Loading data and cleaning}\label{loading-data-and-cleaning}

First, I filtered and cleaned the title data. I filtered the data to
only inclue movies. I cleaned the data to remove column formatting
issues. As some of the columns had two values per observation
(e.g.~multiple genres), I pivoted longer to allow for genre and country
based analysis.

\section{Types of movies by country}\label{types-of-movies-by-country}

I first identify the identify the top countries. The five countries
which are the most prolific producers are the United States, India,
France, and Canada.

\begin{longtable}[]{@{}llr@{}}
\caption{Top 5 Genres by Country}\tabularnewline
\toprule\noalign{}
production\_countries & genres & n \\
\midrule\noalign{}
\endfirsthead
\toprule\noalign{}
production\_countries & genres & n \\
\midrule\noalign{}
\endhead
\bottomrule\noalign{}
\endlastfoot
US & comedy & 630 \\
US & drama & 618 \\
US & documentation & 365 \\
US & thriller & 353 \\
US & action & 251 \\
IN & drama & 448 \\
IN & comedy & 252 \\
IN & romance & 192 \\
IN & thriller & 174 \\
IN & action & 149 \\
GB & drama & 116 \\
GB & thriller & 80 \\
GB & comedy & 77 \\
GB & documentation & 59 \\
GB & action & 49 \\
FR & drama & 113 \\
FR & european & 100 \\
FR & comedy & 51 \\
FR & thriller & 51 \\
FR & crime & 36 \\
CA & drama & 63 \\
CA & comedy & 56 \\
CA & thriller & 40 \\
CA & animation & 31 \\
CA & family & 31 \\
\end{longtable}

The results from analysis the top genres by countries illustrate
heterogeneity in preferences. Comedy and drama dominate in the United
States, India, and Canada. Documentation films are strong in the United
States, and the UK.

\begin{figure}[H]

{\centering \includegraphics{22921893Question4_files/figure-latex/unnamed-chunk-4-1} 

}

\caption{Top Generes by Country.\label{fig:country-genre}}\label{fig:unnamed-chunk-4}
\end{figure}

\section{Distribution of ratings by
genre}\label{distribution-of-ratings-by-genre}

Here, I visualise the distribution of ratings by genre. History,
documentation, war and sport films are more commonly well-received.
Horror, family, sci-fi, and action lie at the bottom range of the rating
distribution. However, most genres cluster around a mean score of
between 6 and 7, and are fairly widely spread. This indicates genre
alone is a weak predictor of ratings.

\begin{figure}[H]

{\centering \includegraphics{22921893Question4_files/figure-latex/unnamed-chunk-5-1} 

}

\caption{IMDb Score Distribution by Genre.\label{fig:genreridge}}\label{fig:unnamed-chunk-5}
\end{figure}

\section{Distribution of runtime by
genre}\label{distribution-of-runtime-by-genre}

Here, I visualise the distribution of runtime by genre. Similar to the
ratings distribution, most genres are clustered over a similar mean
(just below 100 minutes) with fairly wide spreads.

\begin{figure}[H]

{\centering \includegraphics{22921893Question4_files/figure-latex/unnamed-chunk-6-1} 

}

\caption{Movie Runtime Distribution by Genre.\label{fig:runtimeviolin}}\label{fig:unnamed-chunk-6}
\end{figure}

\section{Text analysis}\label{text-analysis}

First, I clean and filter the dataset. To avoid double counting words
from movies with multiple genres, the data is reduced to one row per
movie per country before tokenising descriptions.

Second, I look for the top words by country. Across most countries,
``life'', ``family'', and ``world'' are common in movie descriptions.
This indicates that a strong drama/relationship focus is appealing. The
UK stands out with location-specific words like ``london'' and
``british''. Similarly, France's results lean toward ``french'' and
``paris''. This suggesting more nationally-rooted storytelling is
favoured in these countries. In contrast, Canada and the United states
have more broad themes, such as ``family'', ``life'' and ``home''.

\begin{figure}[H]

{\centering \includegraphics{22921893Question4_files/figure-latex/unnamed-chunk-9-1} 

}

\caption{Top 10 Description Words by Country.\label{fig:wordplot}}\label{fig:unnamed-chunk-9}
\end{figure}

South Africa's most frequent words have a slightly different tilt than
the world's top producers. Notably, frequent words include ``south'' and
``africa'', indicating the important of domestic storytelling in the
country. ``Apartheid'' is also amongst the most frequent, illustrating
the importance of the countries history. Similar to the patterns
observed in the rest of the world, words such as ``love'', ``life'' and
``family'' are also common in South Africa.

\begin{figure}[H]

{\centering \includegraphics{22921893Question4_files/figure-latex/unnamed-chunk-10-1} 

}

\caption{Top Description Words South Africa.\label{fig:rsalollipop}}\label{fig:unnamed-chunk-10}
\end{figure}

\bibliography{Tex/ref}





\end{document}
